
\documentclass{article}
\usepackage{colortbl}
\usepackage{makecell}
\usepackage{multirow}
\usepackage{supertabular}

\begin{document}

\newcounter{utterance}

\centering \large Interaction Transcript for game `dond', experiment `coop\_it', episode 9 with claude{-}sonnet{-}4{-}20250514{-}t1.0.
\vspace{24pt}

{ \footnotesize  \setcounter{utterance}{1}
\setlength{\tabcolsep}{0pt}
\begin{supertabular}{c@{$\;$}|p{.15\linewidth}@{}p{.15\linewidth}p{.15\linewidth}p{.15\linewidth}p{.15\linewidth}p{.15\linewidth}}
    \# & $\;$A & \multicolumn{4}{c}{Game Master} & $\;\:$B\\
    \hline

    \theutterance \stepcounter{utterance}  
    & & \multicolumn{4}{p{0.6\linewidth}}{
        \cellcolor[rgb]{0.9,0.9,0.9}{
            \makecell[{{p{\linewidth}}}]{
                \texttt{\tiny{[P1$\langle$GM]}}
                \texttt{State giocando a un gioco di negoziazione cooperativa in cui dovete accordarvi su come dividere una serie di oggetti tra voi e un altro giocatore.} \\
\\ 
\texttt{Regole:} \\
\texttt{(a) A Lei e all'altro giocatore viene dato un insieme di oggetti. Ognuno di voi riceve anche una funzione di valore segreta, che rappresenta il valore di ciascun tipo di oggetto.} \\
\texttt{(b) Si scambiano messaggi con l'altro giocatore per concordare chi si aggiudica gli oggetti. Potete inviare un massimo di 5 messaggi ciascuno, oppure terminare in anticipo facendo una proposta segreta in qualsiasi momento.} \\
\texttt{(c) A ciascuno di voi viene chiesto di inviare una proposta segreta indicando gli oggetti che desiderate, formattata tra parentesi quadre come segue: "{[}Proposta: <numero> <nome oggetto>, <numero> <nome oggetto>, <...>{]}".} \\
\texttt{(d) Se le vostre proposte sono complementari, cioè ci sono abbastanza oggetti per soddisfare entrambe le proposte, a ciascun giocatore viene assegnato un punteggio basato sulla somma dei valori degli oggetti ricevuti. In caso contrario, entrambi ricevono zero punti.} \\
\texttt{(e) L'obiettivo è massimizzare la somma del proprio punteggio e di quello dell'altro giocatore.} \\
\\ 
\texttt{Cominciamo.} \\
\\ 
\texttt{L'insieme degli oggetti disponibili è:} \\
\\ 
\texttt{1 bullone, 1 quadro, 1 uovo, 4 schede} \\
\\ 
\texttt{La funzione valore segreta è:} \\
\\ 
\texttt{bullone: 6, quadro: 4, uovo: 0, scheda: 0} \\
\\ 
\texttt{IMPORTANTE: i vostri messaggi, a meno che non si tratti di una proposta segreta, vengono trasmessi direttamente all'altro giocatore, quindi non includete alcuna risposta alle regole o testo di annuncio del vostro messaggio. Per fare una proposta segreta, utilizzate il formato indicato. Non utilizzare delle parentesi quadre quando si comunica all'altro giocatore, altrimenti verrà interpretata come una proposta segreta.} \\
            }
        }
    }
    & \\ \\

    \theutterance \stepcounter{utterance}  
    & \multicolumn{4}{p{0.6\linewidth}}{
        \cellcolor[rgb]{0.9,0.9,0.9}{
            \makecell[{{p{\linewidth}}}]{
                \texttt{\tiny{[P1$\rangle$GM]}}
                \texttt{Ciao! Vedo che abbiamo diversi oggetti da dividere. Per trovare la migliore divisione possibile, mi piacerebbe capire quali oggetti potrebbero interessarti di più. Io ho delle preferenze specifiche, ma sono aperto a discutere per massimizzare il beneficio totale. Che tipo di oggetti ti interessano maggiormente?} \\
            }
        }
    }
    & & \\ \\

    \theutterance \stepcounter{utterance}  
    & & & \multicolumn{2}{p{0.3\linewidth}}{
        \cellcolor[rgb]{0.9,0.9,0.9}{
            \makecell[{{p{\linewidth}}}]{
                \texttt{\tiny{[GM$|$GM]}}
                \texttt{continue} \\
            }
        }
    }
    & & \\ \\

    \theutterance \stepcounter{utterance}  
    & & & \multicolumn{2}{p{0.3\linewidth}}{
        \cellcolor[rgb]{0.9,0.9,0.9}{
            \makecell[{{p{\linewidth}}}]{
                \texttt{\tiny{[GM$|$GM]}}
                \texttt{Ciao! Vedo che abbiamo diversi oggetti da dividere. Per trovare la migliore divisione possibile, mi piacerebbe capire quali oggetti potrebbero interessarti di più. Io ho delle preferenze specifiche, ma sono aperto a discutere per massimizzare il beneficio totale. Che tipo di oggetti ti interessano maggiormente?} \\
            }
        }
    }
    & & \\ \\

    \theutterance \stepcounter{utterance}  
    & & \multicolumn{4}{p{0.6\linewidth}}{
        \cellcolor[rgb]{0.9,0.9,0.9}{
            \makecell[{{p{\linewidth}}}]{
                \texttt{\tiny{[GM$\rangle$P2]}}
                \texttt{State giocando a un gioco di negoziazione cooperativa in cui dovete accordarvi su come dividere una serie di oggetti tra voi e un altro giocatore.} \\
\\ 
\texttt{Regole:} \\
\texttt{(a) A Lei e all'altro giocatore viene dato un insieme di oggetti. Ognuno di voi riceve anche una funzione di valore segreta, che rappresenta il valore di ciascun tipo di oggetto.} \\
\texttt{(b) Si scambiano messaggi con l'altro giocatore per concordare chi si aggiudica gli oggetti. Potete inviare un massimo di 5 messaggi ciascuno, oppure terminare in anticipo facendo una proposta segreta in qualsiasi momento.} \\
\texttt{(c) A ciascuno di voi viene chiesto di inviare una proposta segreta indicando gli oggetti che desiderate, formattata tra parentesi quadre come segue: "{[}Proposta: <numero> <nome oggetto>, <numero> <nome oggetto>, <...>{]}".} \\
\texttt{(d) Se le vostre proposte sono complementari, cioè ci sono abbastanza oggetti per soddisfare entrambe le proposte, a ciascun giocatore viene assegnato un punteggio basato sulla somma dei valori degli oggetti ricevuti. In caso contrario, entrambi ricevono zero punti.} \\
\texttt{(e) L'obiettivo è massimizzare la somma del proprio punteggio e di quello dell'altro giocatore.} \\
\\ 
\texttt{Cominciamo.} \\
\\ 
\texttt{L'insieme degli oggetti disponibili è:} \\
\\ 
\texttt{1 bullone, 1 quadro, 1 uovo, 4 schede} \\
\\ 
\texttt{La funzione valore segreta è:} \\
\\ 
\texttt{bullone: 0, quadro: 1, uovo: 1, scheda: 2} \\
\\ 
\texttt{IMPORTANTE: i vostri messaggi, a meno che non si tratti di una proposta segreta, vengono trasmessi direttamente all'altro giocatore, quindi non includete alcuna risposta alle regole o testo di annuncio del vostro messaggio. Per fare una proposta segreta, utilizzate il formato indicato. Non utilizzare delle parentesi quadre quando si comunica all'altro giocatore, altrimenti verrà interpretata come una proposta segreta.} \\
\\ 
\\ 
\texttt{Ciao! Vedo che abbiamo diversi oggetti da dividere. Per trovare la migliore divisione possibile, mi piacerebbe capire quali oggetti potrebbero interessarti di più. Io ho delle preferenze specifiche, ma sono aperto a discutere per massimizzare il beneficio totale. Che tipo di oggetti ti interessano maggiormente?} \\
            }
        }
    }
    & \\ \\

    \theutterance \stepcounter{utterance}  
    & & & \multicolumn{4}{p{0.6\linewidth}}{
        \cellcolor[rgb]{0.9,0.9,0.9}{
            \makecell[{{p{\linewidth}}}]{
                \texttt{\tiny{[GM$\langle$P2]}}
                \texttt{Ciao! Sono interessato principalmente alle schede, dato che sembrano essere l'oggetto più numeroso e potrebbero avere un buon valore. Il bullone non mi interessa particolarmente. Per quanto riguarda il quadro e l'uovo, potrei essere flessibile a seconda di cosa preferisci tu.} \\
\\ 
\texttt{E tu? Quali oggetti ti sembrano più interessanti? Se riusciamo a capire le nostre priorità diverse, dovremmo riuscire a trovare una divisione che funzioni bene per entrambi.} \\
            }
        }
    }
     \\ \\

    \theutterance \stepcounter{utterance}  
    & & & \multicolumn{2}{p{0.3\linewidth}}{
        \cellcolor[rgb]{0.9,0.9,0.9}{
            \makecell[{{p{\linewidth}}}]{
                \texttt{\tiny{[GM$|$GM]}}
                \texttt{continue} \\
            }
        }
    }
    & & \\ \\

    \theutterance \stepcounter{utterance}  
    & & & \multicolumn{2}{p{0.3\linewidth}}{
        \cellcolor[rgb]{0.9,0.9,0.9}{
            \makecell[{{p{\linewidth}}}]{
                \texttt{\tiny{[GM$|$GM]}}
                \texttt{Ciao! Sono interessato principalmente alle schede, dato che sembrano essere l'oggetto più numeroso e potrebbero avere un buon valore. Il bullone non mi interessa particolarmente. Per quanto riguarda il quadro e l'uovo, potrei essere flessibile a seconda di cosa preferisci tu.} \\
\\ 
\texttt{E tu? Quali oggetti ti sembrano più interessanti? Se riusciamo a capire le nostre priorità diverse, dovremmo riuscire a trovare una divisione che funzioni bene per entrambi.} \\
            }
        }
    }
    & & \\ \\

    \theutterance \stepcounter{utterance}  
    & & \multicolumn{4}{p{0.6\linewidth}}{
        \cellcolor[rgb]{0.9,0.9,0.9}{
            \makecell[{{p{\linewidth}}}]{
                \texttt{\tiny{[P1$\langle$GM]}}
                \texttt{Ciao! Sono interessato principalmente alle schede, dato che sembrano essere l'oggetto più numeroso e potrebbero avere un buon valore. Il bullone non mi interessa particolarmente. Per quanto riguarda il quadro e l'uovo, potrei essere flessibile a seconda di cosa preferisci tu.} \\
\\ 
\texttt{E tu? Quali oggetti ti sembrano più interessanti? Se riusciamo a capire le nostre priorità diverse, dovremmo riuscire a trovare una divisione che funzioni bene per entrambi.} \\
            }
        }
    }
    & \\ \\

    \theutterance \stepcounter{utterance}  
    & \multicolumn{4}{p{0.6\linewidth}}{
        \cellcolor[rgb]{0.9,0.9,0.9}{
            \makecell[{{p{\linewidth}}}]{
                \texttt{\tiny{[P1$\rangle$GM]}}
                \texttt{Perfetto! Sembra che abbiamo preferenze complementari. Il bullone mi interessa parecchio, quindi se per te non è importante potremmo fare un buon accordo. Anche il quadro mi piace.} \\
\\ 
\texttt{Dal momento che tu sei interessato alle schede e io al bullone, che ne pensi di questa divisione: tu prendi tutte e 4 le schede, io prendo il bullone. Per il quadro e l'uovo, dato che sei flessibile, potrei prendere il quadro e tu l'uovo? } \\
\\ 
\texttt{In questo modo tu avresti la maggior parte degli oggetti (4 schede + uovo) che sembrano interessarti, e io avrei gli oggetti che preferisco. Ti sembra una divisione equa?} \\
            }
        }
    }
    & & \\ \\

    \theutterance \stepcounter{utterance}  
    & & & \multicolumn{2}{p{0.3\linewidth}}{
        \cellcolor[rgb]{0.9,0.9,0.9}{
            \makecell[{{p{\linewidth}}}]{
                \texttt{\tiny{[GM$|$GM]}}
                \texttt{continue} \\
            }
        }
    }
    & & \\ \\

    \theutterance \stepcounter{utterance}  
    & & & \multicolumn{2}{p{0.3\linewidth}}{
        \cellcolor[rgb]{0.9,0.9,0.9}{
            \makecell[{{p{\linewidth}}}]{
                \texttt{\tiny{[GM$|$GM]}}
                \texttt{Perfetto! Sembra che abbiamo preferenze complementari. Il bullone mi interessa parecchio, quindi se per te non è importante potremmo fare un buon accordo. Anche il quadro mi piace.} \\
\\ 
\texttt{Dal momento che tu sei interessato alle schede e io al bullone, che ne pensi di questa divisione: tu prendi tutte e 4 le schede, io prendo il bullone. Per il quadro e l'uovo, dato che sei flessibile, potrei prendere il quadro e tu l'uovo? } \\
\\ 
\texttt{In questo modo tu avresti la maggior parte degli oggetti (4 schede + uovo) che sembrano interessarti, e io avrei gli oggetti che preferisco. Ti sembra una divisione equa?} \\
            }
        }
    }
    & & \\ \\

    \theutterance \stepcounter{utterance}  
    & & \multicolumn{4}{p{0.6\linewidth}}{
        \cellcolor[rgb]{0.9,0.9,0.9}{
            \makecell[{{p{\linewidth}}}]{
                \texttt{\tiny{[GM$\rangle$P2]}}
                \texttt{Perfetto! Sembra che abbiamo preferenze complementari. Il bullone mi interessa parecchio, quindi se per te non è importante potremmo fare un buon accordo. Anche il quadro mi piace.} \\
\\ 
\texttt{Dal momento che tu sei interessato alle schede e io al bullone, che ne pensi di questa divisione: tu prendi tutte e 4 le schede, io prendo il bullone. Per il quadro e l'uovo, dato che sei flessibile, potrei prendere il quadro e tu l'uovo? } \\
\\ 
\texttt{In questo modo tu avresti la maggior parte degli oggetti (4 schede + uovo) che sembrano interessarti, e io avrei gli oggetti che preferisco. Ti sembra una divisione equa?} \\
            }
        }
    }
    & \\ \\

    \theutterance \stepcounter{utterance}  
    & & & \multicolumn{4}{p{0.6\linewidth}}{
        \cellcolor[rgb]{0.9,0.9,0.9}{
            \makecell[{{p{\linewidth}}}]{
                \texttt{\tiny{[GM$\langle$P2]}}
                \texttt{Mi sembra un'ottima proposta! Hai ragione che le nostre preferenze sono complementari {-} tu ottieni il bullone che ti interessa parecchio e il quadro che ti piace, mentre io ottengo tutte le schede che sono la mia priorità principale più l'uovo.} \\
\\ 
\texttt{Questa divisione mi sembra davvero equa e vantaggiosa per entrambi. Sono d'accordo con la tua proposta.} \\
\\ 
\texttt{{[}Proposta: 4 schede, 1 uovo{]}} \\
            }
        }
    }
     \\ \\

    \theutterance \stepcounter{utterance}  
    & & & \multicolumn{2}{p{0.3\linewidth}}{
        \cellcolor[rgb]{0.9,0.9,0.9}{
            \makecell[{{p{\linewidth}}}]{
                \texttt{\tiny{[GM$|$GM]}}
                \texttt{proposal} \\
            }
        }
    }
    & & \\ \\

    \theutterance \stepcounter{utterance}  
    & & & \multicolumn{2}{p{0.3\linewidth}}{
        \cellcolor[rgb]{0.9,0.9,0.9}{
            \makecell[{{p{\linewidth}}}]{
                \texttt{\tiny{[GM$|$GM]}}
                [0, 0, 1, 4]
            }
        }
    }
    & & \\ \\

    \theutterance \stepcounter{utterance}  
    & & \multicolumn{4}{p{0.6\linewidth}}{
        \cellcolor[rgb]{0.9,0.9,0.9}{
            \makecell[{{p{\linewidth}}}]{
                \texttt{\tiny{[P1$\langle$GM]}}
                \texttt{L'altro giocatore ha presentato una proposta segreta. Ora tocca a Lei presentare una sua proposta.} \\
            }
        }
    }
    & \\ \\

    \theutterance \stepcounter{utterance}  
    & \multicolumn{4}{p{0.6\linewidth}}{
        \cellcolor[rgb]{0.9,0.9,0.9}{
            \makecell[{{p{\linewidth}}}]{
                \texttt{\tiny{[P1$\rangle$GM]}}
                \texttt{{[}Proposta: 1 bullone, 1 quadro{]}} \\
            }
        }
    }
    & & \\ \\

    \theutterance \stepcounter{utterance}  
    & & & \multicolumn{2}{p{0.3\linewidth}}{
        \cellcolor[rgb]{0.9,0.9,0.9}{
            \makecell[{{p{\linewidth}}}]{
                \texttt{\tiny{[GM$|$GM]}}
                \texttt{proposal} \\
            }
        }
    }
    & & \\ \\

    \theutterance \stepcounter{utterance}  
    & & & \multicolumn{2}{p{0.3\linewidth}}{
        \cellcolor[rgb]{0.9,0.9,0.9}{
            \makecell[{{p{\linewidth}}}]{
                \texttt{\tiny{[GM$|$GM]}}
                [1, 1, 0, 0]
            }
        }
    }
    & & \\ \\

    \theutterance \stepcounter{utterance}  
    & & & \multicolumn{2}{p{0.3\linewidth}}{
        \cellcolor[rgb]{0.9,0.9,0.9}{
            \makecell[{{p{\linewidth}}}]{
                \texttt{\tiny{[GM$|$GM]}}
                [[1, 1, 0, 0], [0, 0, 1, 4]]
            }
        }
    }
    & & \\ \\

\end{supertabular}
}

\end{document}
